\clearpage
\thispagestyle{plain}
\phantomsection
\addcontentsline{toc}{chapter}{Abstract}

\begin{center}
{\bfseries\Large ABSTRACT\par}

\end{center}

\vspace{0.8cm}

\noindent
FireWatch is a drone-ready wildfire risk prediction and fire/smoke detection
prototype developed for the Karabük region in Türkiye. The project addresses
the need for an integrated monitoring workflow that combines early
weather-based risk awareness with visual evidence from camera sources. The
prediction methodology used May--October seasonal Karabük records from 2012
to 2025. Open-Meteo daily weather aggregates were used as accessible model
inputs, while the dataset included a processed Fire Weather Index
(FWI) target associated with official fire-danger data and used for supervised
learning.

The prediction pipeline was designed as a two-stage model. Stage 1 used a
HistGradientBoostingRegressor to estimate continuous FWI from a locked
34-feature schema. Stage 2 used a RandomForestClassifier as a safety-oriented
high-risk classifier for days where FWI is greater than or equal to 35. For
visual monitoring, a custom YOLOv8 model, best.pt, was used to detect fire
and smoke from live camera or drone-ready sources.

The system was implemented using a FastAPI backend, SQLite persistence,
scheduled risk checks, and a React/TypeScript dashboard. The dashboard
integrates risk decisions, live monitoring feeds, detection alerts,
model/system information, workflow explanation, and map context. The results
show that FireWatch can integrate prediction, detection, persistence, and
dashboard-based decision support into one regional wildfire monitoring
prototype. The system remains an academic prototype and does not claim
official emergency deployment or full autonomous drone operation.

\vspace{0.8cm}

\noindent\textbf{Key Words:} wildfire prediction, fire detection, smoke detection, Fire Weather Index, YOLOv8, drone-ready monitoring, decision-support system

\clearpage
